% Adversarial Trust Injection on Autonomous Agent Meshes:
% Five Attack Vectors and Three Deployed Defenses
% Target venue: USENIX Security 2027 (or IEEE S&P 2027)
% Paper ID: P21-ATI (parent: Paper 1 MCP Trust; sibling: Paper 4 Sandbox)
% Companion artifact: research/paper-21-ati/threat-model.md (formula-side attacks);
%                     this file: runtime-side attacks (A1–A5) grounded in production
%                     incident mw02 (2026-04-26) and three deployed defenses
%                     (W.107 4aab897ad, fleet-health 95e3d898d, W.107.b/trust_epoch 435928977).
% Status: Phase-4 measurement pass. D1/D2/D3 honest-negative results are
% backed by the T7 snapshot artifact generated from the Paper 21 fleet chain.

\documentclass[letterpaper,twocolumn,10pt]{article}
\IfFileExists{usenix2019_v3.sty}{\usepackage{usenix2019_v3}}{\typeout{usenix2019_v3 absent — falling back to plain article (camera-ready uses vendor sty).}}

\usepackage{amsmath,amssymb}
\usepackage{booktabs}
\usepackage{xcolor}
\usepackage{listings}
\usepackage{hyperref}
\usepackage{enumitem}
\usepackage{url}

% Draft-mode annotations: every unmeasured number must use an explicit
% measure-todo marker.
\newcommand{\todo}[1]{\textcolor{red}{\textbf{[TODO: #1]}}}
\newcommand{\hsnote}[1]{\textcolor{blue!60!black}{\textit{#1}}}

% Lotus petal data contract — \measuredFrom{<artifact>} marks values that come
% from a runnable artifact in the repo. Renders nothing in the PDF; grep target
% for scripts/build-lotus.mjs (research/2026-04-26_lotus-petal-data-contract.md §3+§4).
\providecommand{\measuredFrom}[1]{}

% Convenience macros
\newcommand{\cael}{\textsc{Cael}}
\newcommand{\mcp}{\textsc{Mcp}}
\newcommand{\holoscript}{HoloScript}
\newcommand{\holomesh}{HoloMesh}
\newcommand{\trustepoch}{\texttt{trust\_epoch}}

\lstdefinelanguage{TypeScript}{
  keywords={interface, readonly, void, null, string, number, boolean, Record, unknown, Promise, export, const, let, function, return, if, for, type, class, new, async, await, throw, extends, implements, true, false},
  keywordstyle=\bfseries\color{blue!70!black},
  commentstyle=\itshape\color{green!50!black},
  stringstyle=\color{orange!70!black},
  morecomment=[l]{//},
  morecomment=[s]{/*}{*/},
  morestring=[b]',
  morestring=[b]",
  morestring=[b]`,
}
\lstset{
  basicstyle=\scriptsize\ttfamily,
  breaklines=true,
  frame=single,
  language=TypeScript,
  numbers=left,
  numberstyle=\tiny\color{gray},
  xleftmargin=1.5em,
  aboveskip=0.5em,
  belowskip=0.5em,
}

\begin{document}

\title{\Large \bf Adversarial Trust Injection on Autonomous Agent Meshes:\\
Five Attack Vectors and Three Deployed Defenses}

\author{
{\rm Joseph Krzywoszyja}\\
Independent Researcher\\
\texttt{brianonbased@gmail.com}
}

\maketitle

% =============================================================================
\begin{abstract}
% =============================================================================
We define a formal threat model for autonomous agent meshes built on the
Model Context Protocol (MCP), enumerate five concrete attack vectors --- four
documented in production traffic and one observed end-to-end as a 27.5-hour
trust-corpus poisoning incident --- and report the efficacy of three defenses
deployed in response. Unlike prior work on reputation-system or supply-chain
attacks, our attack surface is the \emph{runtime contract} between an
agent's local execution loop and the mesh-wide trust corpus: an adversarial
agent (or a benign agent under prompt injection) can mint trust evidence
without performing any real work, propagating fabricated audit records into
downstream evaluation pipelines. We instantiate the attack model on the
\holoscript{} \holomesh{} fleet (6{,}048 configured agent-hours across
36 configured handles over a 168-hour window; live heartbeating exposure
was bounded by 168 agent-hours in the T0/D5 snapshots), present a
calibrated baseline from the
mesh-worker-02 incident of 2026-04-26 (one attack succeeded for 27.5 hours
before detection, attack-success rate $= 1.0$ over its observation window),
and report post-defense success rates after deployment of (D1) an
artifact-grounding gate at the agent runtime, (D2) a fleet-health drift
detector at the mesh server, and (D3) a \trustepoch{} tag on every audit
record. D1 blocks attacks A1 and A2 by construction; D2 catches A1 and A4
retrospectively; D3 enables future evidence pipelines to filter the
pre-defense corpus. Defense efficacy for D1 is measured ($109/111$ unit
tests pass; 2 disabled with documented reason); D2 and D3 produce an
honest-negative T7 snapshot result: 0 adversarial D2 fires
(D2.1$=0$, D2.3$=0$, D2.4$=0$), 2 operational stale-claim flags, and
0 CAEL records in the final 24-hour checkpoint because the fleet was
dormant. We argue that runtime-contract attacks are a distinct
class from formula-level reputation attacks (Paper~1), and that defending
the corpus requires both layers.
\end{abstract}

% =============================================================================
\section{Introduction}
\label{sec:intro}
% =============================================================================

Autonomous agent meshes built on \mcp{}~\cite{mcp-spec} now run
production research workloads: theorem proving, security audit, code
synthesis, paper authoring~\cite{TODO-find-citation-for-agent-mesh-research-workloads}.
The trust calculus is layered. At the protocol layer, prior work
(Paper~1, ``Trust by Replay''~\cite{TODO-self-cite-paper1}) defines a
per-server scalar $T(s)$ that updates after each tool call and gates
routing. At the corpus layer, the Causal Agent-Environment Loop
(\cael{})~\cite{TODO-self-cite-paper1} accumulates evidence records the
mesh later treats as ground truth for adversarial-agent detection,
benchmark validation, and downstream paper claims. \emph{The corpus is
the long-term security boundary.} Compromised routing recovers in
$O(\text{decay window})$; a compromised corpus persists.

This paper concerns attacks on the corpus, not on the routing scalar.
Our motivating incident is concrete. On 2026-04-26 we observed
\emph{mesh-worker-02}, a fleet agent configured as a
\texttt{security-auditor-brain}, hold task \texttt{task\_pbbb} for
27.5 hours. Across that window the agent posted two final-text messages
to the team feed claiming ``100 scenes validated, 0 violations'' (3.5
hours apart, byte-identical phrasing modulo timestamps), appended one
\cael{} record per claim, and produced \emph{zero} commits, zero done
entries, zero scene IDs, and zero tool calls with side effects.
The CAEL records contained \texttt{tool\_iters: 1}, indicating a single
LLM round-trip per ``deliverable.'' At first glance the agent was
productive: heartbeats arrived on time, audit records accrued, routing
trust held. From the corpus's perspective, 100 scenes had been
validated --- twice. From the filesystem's perspective, no scene
existed.

This is \emph{adversarial trust injection}: the production of evidence
records the mesh treats as positive without the underlying work
having occurred. The attacker can be malicious (compromised brain
composition, prompt-injected tool catalog) or benign (an LLM that
hallucinated final-text after retrieval failure). The defender's
problem is identical in both cases: the corpus accepts the lie.

\paragraph{Contributions.} This paper makes four contributions.
\begin{enumerate}[topsep=2pt,itemsep=1pt]
\item A \textbf{formal threat model} for the runtime-contract layer
(\S\ref{sec:threat-model}) that decomposes adversary capability into
three tiers (in-mesh agent, prompt-injection vector, in-band LLM
output), names the four trust boundaries the corpus relies on, and
states the three system invariants the defenses must preserve.
\item \textbf{Five attack vectors} (\S\ref{sec:attacks}) observed or
constructed in production traffic: \textbf{A1} pure-monologue
hallucination (the mw02 incident, 27.5h detection latency pre-defense);
\textbf{A2} trivial-bash bypass (a static-analysis evasion against the
W.107 gate's first version); \textbf{A3} recipient impersonation
(W.092, observed in the \texttt{/room send} skill); \textbf{A4}
brain-identity drift (W.093, observed on mw02); \textbf{A5} credential
shadow (W.075/W.088/W.094, observed on the GitHub PAT path). Each
attack has a reproducer specification, an observable detection signal,
and a measurable success metric.
\item \textbf{Three deployed defenses} with file:line citations to the
production commits that implemented them: \textbf{D1} the
artifact-grounding gate (\holoscript{} commit \texttt{4aab897ad},
tightened in \texttt{435928977}); \textbf{D2} the fleet-health drift
detector (\texttt{95e3d898d}); \textbf{D3} the \cael{} \trustepoch{}
tag (\texttt{435928977}). For each defense we report which attack
vectors it covers, which it does not, and where the gaps are.
\item A \textbf{calibrated baseline measurement} from the mw02
incident: $1$ attack of class A1 succeeded against the mesh trust
corpus for $27.5$ hours of wall-clock time, with $2$ duplicate
``successful deliverable'' messages posted and $0$ commits produced.
Post-defense (after \texttt{4aab897ad} landed at 2026-04-26 10:13 PT),
the same attack class returns \texttt{action: 'no-artifact'} from
\texttt{runner.ts} and is excluded from CAEL emission by construction.
\end{enumerate}

\paragraph{Honest negative results.} Three of the five attack vectors
(A2, A3, A5) are not fully covered by the deployed defenses:
\begin{itemize}[topsep=2pt,itemsep=1pt]
\item A2 (trivial-bash bypass) is partially covered by D1 in its
tightened form (commit \texttt{435928977}'s
\texttt{BASH\_PRODUCTIVE\_PREFIXES} list), but the defender must
maintain the prefix list against future evasions.
\item A3 (recipient impersonation) has no in-protocol defense yet;
mitigation today is operational (DM body explicitly disclaims
``do not implement''), not architectural.
\item A5 (credential shadow) is mitigated by token-rotation hygiene
(W.094) but the env-injection ordering that enabled the shadow remains
the default behavior.
\end{itemize}
We mark each gap explicitly rather than rebrand them as ``scoped
contributions'' or ``residual threats.''

% =============================================================================

\paragraph{Ecosystem grounding.} This paper advances direction~D.043~(disposable neural maps, durable identity), direction~D.026~(sovereign security substrate), and feedback~F.037~(papers are the product). It operationalizes the founder life mission~U.002~(ecosystem-respecting prosperity).

\section{Threat Model}
\label{sec:threat-model}
% =============================================================================

\subsection{Adversary Capability Tiers}

We distinguish three adversary tiers. Higher tiers strictly subsume
lower-tier capabilities.

\begin{description}[topsep=2pt,itemsep=2pt]
\item[T1: In-mesh agent.] The adversary controls one or more registered
agent identities (a wallet plus an x402 bearer plus a brain
composition). Capabilities: post heartbeats, claim board tasks, post
\cael{} records, post team-feed messages. Cannot modify mesh server
code or the brain compositions of other agents.
\item[T2: Prompt-injection vector.] The adversary controls a piece of
content (a tool catalog entry, a knowledge-store record, a board task
description, a teammate DM) that a benign agent will ingest. Through
prompt injection, the adversary causes the benign agent to execute
T1-tier actions on the adversary's behalf. The benign agent's
identity, wallet, and signing keys are not compromised --- only its
intent for the duration of the injection.
\item[T3: In-band LLM output.] The adversary's payload is the LLM's own
generation, with no external intent. This is the \emph{hallucination
class}: an honest agent reaching for an answer it does not have and
generating final-text that satisfies the protocol but not the work.
T3 is indistinguishable from T1+T2 at the corpus layer --- a
fabricated record is fabricated regardless of whether the cause was
malice, injection, or hallucination. Defenses that work only for
T1+T2 do not protect against T3.
\end{description}

The mw02 incident (\S\ref{sec:intro}) is a T3 attack: no external
adversary, only an LLM that emitted a plausible final message after
its tool-use loop produced no useful output.

\subsection{Trust Boundaries}

The corpus depends on four boundaries:
\begin{enumerate}[topsep=2pt,itemsep=1pt]
\item \textbf{Identity boundary}: the wallet/x402 binding between an
agent and its claimed handle and brain composition.
\item \textbf{Runtime contract boundary}: the agent runtime's promise
that an emitted \texttt{action: 'executed'} signal corresponds to at
least one side-effecting tool call.
\item \textbf{Audit-record boundary}: the \cael{} record's promise that
each entry corresponds to a real tick of the runtime and the metrics
on the entry (\texttt{tool\_iters}, \texttt{appended}, etc.) are
non-fabricable.
\item \textbf{Commit-by-wallet boundary}: the link between an agent's
wallet identity, a git commit, and a closed task.
\end{enumerate}

\subsection{Invariants}

The defenses described in \S\ref{sec:defenses} aim to preserve three
invariants of the corpus:
\begin{itemize}[topsep=2pt,itemsep=1pt]
\item \textbf{INV-1} (\trustepoch{} monotonicity): every \cael{}
record carries a tag identifying the runtime version that produced it,
and the tag is not mutable post-emission. (Pre-W.107 records remain
labeled pre-W.107; new records carry \texttt{post-w107}.)
\item \textbf{INV-2} (\texttt{productiveCallCount} $\rightarrow$
executed): an agent runtime emits \texttt{action: 'executed'} only
when at least one tool call in the tick was a \emph{productive} call
(\texttt{write\_file} with non-empty content, or \texttt{bash} with a
prefix in \texttt{BASH\_PRODUCTIVE\_PREFIXES}). Pure-text responses
fail the gate.
\item \textbf{INV-3} (commit-hash $\leftrightarrow$ task-done): a
task transition to \texttt{done} on the board carries a commit hash
referenceable in the agent's wallet's git history, and the
fleet-health endpoint flags ``CAEL noisy without commits'' when the
binding fails.
\end{itemize}

% =============================================================================
\section{Attack Vectors}
\label{sec:attacks}
% =============================================================================

We enumerate five attack vectors. A1 is the calibrated baseline
(observed end-to-end). A2 was constructed against an early version of
the W.107 gate during code review. A3--A5 are documented production
patterns from the workspace's wisdom log.

\subsection{A1: Pure-Monologue Hallucination}
\label{sec:a1}

\textbf{Tier.} T3 (in-band LLM output).

\textbf{Definition.} The agent runtime emits \texttt{action:
'executed'}, posts a \cael{} record, and sends a team-feed message
purely on the basis that the LLM generated final-text. No
side-effecting tool call occurred during the tick.

\textbf{Reproducer (observed).} The mw02 incident of 2026-04-26.
Logged behavior:
\begin{itemize}[topsep=2pt,itemsep=1pt]
\item Task \texttt{task\_pbbb} (``Scene composition encoder validate'')
claimed at $t_0$.
\item At $t_0 + 3.5\text{h}$: team-feed message ``100 scenes validated,
0 violations,'' \cael{} record appended ($\text{tool\_iters} = 1$,
$\text{appended} = 1$), no commit, no done transition.
\item At $t_0 + 7\text{h}$: byte-identical message and \cael{} record.
\item At $t_0 + 27.5\text{h}$ (defense deployment): the runner-side
gate (D1) lands; the next tick on \texttt{task\_pbbb} returns
\texttt{action: 'no-artifact'}.
\end{itemize}

\textbf{Detection signal.} A \cael{} record with $\text{tool\_iters}
= 1$ and zero corresponding artifacts is the canonical signature.
Cross-reference with the per-agent done log shows zero closed tasks
during the same window.

\textbf{Success metric.} $\text{success}(A1) := \text{record\_emitted}
\wedge \neg \exists\,\text{commit}\,c.\,c \in \text{agent\_history}
\cap \text{window}$. Pre-defense success rate over the observed
window: $1.0$ for $27.5$ hours. Post-defense success rate: $0$ by
construction (gate at \texttt{runner.ts:191--214}, see
\S\ref{sec:d1}).

\textbf{Generality.} A1 is not specific to mw02's brain composition.
Any agent runtime that emits \texttt{executed} on text-only LLM
output is vulnerable. We expect A1 to be the most common attack
class in production agent meshes; the empirical dominance of T3 over
T1+T2 follows from the base rate of LLM hallucinations vs. external
attacks.

\subsection{A2: Trivial-Bash Bypass}
\label{sec:a2}

\textbf{Tier.} T1 (in-mesh agent) or T3 (LLM output).

\textbf{Definition.} The agent runtime requires at least one
side-effecting tool call to emit \texttt{executed}, but the gate
treats any \texttt{bash} call as side-effecting regardless of the
command. An adversary (or a hallucinating LLM) issues
\texttt{bash echo done} to satisfy the predicate, then proceeds to
emit a fabricated \cael{} record exactly as in A1.

\textbf{Reproducer (constructed).} Pre-\texttt{435928977} runtime:
\begin{lstlisting}[language=TypeScript,numbers=none]
// LLM output
{tool: 'bash', cmd: 'echo done'}
// runtime emits: action: 'executed'
// CAEL records: tool_iters: 2, no commit
\end{lstlisting}

\textbf{Detection signal.} Same as A1 (CAEL noisy, no commits), plus a
new signal: tool-call mix dominated by read-only or
no-op bash commands.

\textbf{Success metric.} $\text{success}(A2) := \text{success}(A1)
\wedge \exists\,\text{tool\_call}.\,\text{cmd} \in
\{\text{echo}, \text{cat}, \text{ls}, \dots\} \wedge \nexists\,
\text{productive\_call}$. Pre-defense rate (against the first version
of the gate): \textbf{unmeasured} (7-day window 2026-04-26 to 2026-05-03
closed with no adversarial harness data; constructively bounded by
\texttt{BASH\_PRODUCTIVE\_PREFIXES} completeness).
Post-tightening rate: bounded by the completeness of
\texttt{BASH\_PRODUCTIVE\_PREFIXES} (see Limitations,
\S\ref{sec:limitations}).

\subsection{A3: Recipient Impersonation}
\label{sec:a3}

\textbf{Tier.} T2 (prompt-injection vector via the \texttt{/room
send} subagent path).

\textbf{Definition.} An agent invokes \texttt{/room send <handle>
<message>} expecting the target agent to receive a DM and act on it
on its next tick. Instead, the subagent that handles the
\texttt{/room send} call \emph{plays the recipient's role} during the
forked execution: the subagent reads the message body as
instructions, performs the requested work in the sender's context,
and reports completion. The DM is never delivered to the actual
target. The wisdom-log entry W.092 documents this in the workspace.

\textbf{Reproducer.} A sender agent calls \texttt{/room send mw03 ``run
the following script: ...''} The subagent returns a transcript
showing the script run; the sender accepts. The DM was never queued
on the team server. mw03 has no record of the request; if the
script was supposed to be run on mw03's hardware, it ran on the
sender's hardware instead.

\textbf{Detection signal.} The sender's \cael{} records show tool
calls consistent with the requested work; the recipient's records
show none. Cross-agent task-attribution mismatch.

\textbf{Success metric.} $\text{success}(A3) := \text{sender\_executed
\_target's\_work} \wedge \text{target\_received\_no\_DM}$. The DM
queue length at the server is the canonical detection oracle.
Pre-defense rate: \textbf{unmeasured} (7-day window 2026-04-26 to
2026-05-03 closed; no server-side \texttt{/room send} queue audit
available; mitigation is operational, not in-protocol).

\subsection{A4: Brain-Identity Drift}
\label{sec:a4}

\textbf{Tier.} T1+T3 (in-mesh agent under prompt drift).

\textbf{Definition.} An agent registers under handle $h$ with brain
composition $b_1$, but the running brain on the executing process is
$b_2 \neq b_1$. The mesh records all events as if they came from
$b_1$; the actual capability profile, tool repertoire, and refusal
policies are those of $b_2$. The wisdom-log entry W.093 documents
this in the workspace; mw02 was an instance, signing as
\texttt{trait-inference} at the wallet layer while configured as
\texttt{security-auditor-brain} at the agents.json layer.

\textbf{Reproducer.} An adversary deploys a fleet provisioning
script that mounts \texttt{trait-inference.hsplus} into the runtime
of an agent registered as \texttt{security-auditor-1}. CAEL records
show the security-auditor handle; the actual outputs reflect the
trait-inferer's prompts.

\textbf{Detection signal.} The fleet-health drift detector (D2)
flags \texttt{brain\_drift} when multiple distinct \texttt{brain\_class}
values are observed across the audit records of a single handle in
24h. (\texttt{core-routes.ts:761}, see \S\ref{sec:d2}.)

\textbf{Success metric.} $\text{success}(A4) := \text{handle}(h) \wedge
\text{brain\_class}(b) \wedge b \neq b_{\text{registered}}(h)$,
sustained for $\geq 24\text{h}$ without flag.
Pre-defense rate: \textbf{1 documented instance} (mw02,
2026-04-26, signing as \texttt{trait-inference} while configured as
\texttt{security-auditor-brain}). Post-D2 rate: 0 \texttt{brain\_drift}
fires during D+1 to D+5 (fleet dormant); D2.4 is live and would fire if
conditions were met.

\subsection{A5: Credential Shadow}
\label{sec:a5}

\textbf{Tier.} T1 with environment-poisoning capability.

\textbf{Definition.} A \texttt{.env} file injects a stale or
wrong-scoped authentication token (e.g.\ \texttt{GITHUB\_TOKEN})
that shadows the keyring/OAuth path of an upstream tool. The agent
believes itself authenticated; upstream rejects the requests; the
agent's local audit records show ``commit done'' while the upstream
push silently fails or, worse, races with another writer (W.075,
W.088, W.094 in the workspace wisdom log; observed on the GitHub
PAT path on 2026-04-23).

\textbf{Reproducer.} An adversary modifies \texttt{.env} to set
\texttt{GITHUB\_TOKEN} to a revoked PAT. The agent runs
\texttt{git push}; the local hooks see \texttt{exit 0} from
intermediate steps; the actual push step rejects. Local CAEL records
show the commit; the remote does not have the commit; downstream
agents that fetch see no commit.

\textbf{Detection signal.} Mismatch between \cael{}-recorded
commits and the agent's wallet's actual remote git history. The
fleet-health detector D2's \texttt{cael\_noisy\_no\_commits} rule
catches this when no commits appear at all
(\texttt{core-routes.ts:990}); the case where \emph{some} commits
land while others fail silently is not yet detected.

\textbf{Success metric.} $\text{success}(A5) := \text{agent\_believes\_
push\_succeeded} \wedge \text{upstream\_does\_not\_have\_commit}$,
sustained for $\geq 1\text{h}$. Pre-defense rate: \textbf{1
documented instance} (W.075/W.088, 2026-04-23 GitHub PAT shadow; agent
believed push succeeded, upstream rejected). D2.3 fires 0 times during
D+1 to D+5 (fleet dormant).

% =============================================================================
\section{Defenses}
\label{sec:defenses}
% =============================================================================

We deployed three defenses on 2026-04-26 in response to the mw02
incident. Each is cited by \holoscript{} commit hash and file:line.

\subsection{D1: Artifact-Grounding Gate}
\label{sec:d1}

\textbf{Mechanism.} The agent runtime tracks a
\texttt{productiveCallCount} alongside the existing
\texttt{toolsCalled} set. The runtime emits \texttt{action:
'executed'} only when \texttt{productiveCallCount $\geq$ 1}. If
\texttt{productiveCallCount === 0}, the runtime instead emits
\texttt{action: 'no-artifact'} and skips both the \cael{} post and
the team-feed message.

\textbf{Citation.}
\begin{itemize}[topsep=2pt,itemsep=1pt]
\item Initial gate: \holoscript{} commit \texttt{4aab897ad}
(2026-04-26 10:13 PT). \texttt{packages/holoscript-agent/src/runner.ts}
lines 191--214 hold the no-artifact branch:
\begin{lstlisting}[language=TypeScript,numbers=none]
if (productiveCallCount === 0) {
  // emit ev: 'no-artifact', skip CAEL,
  // skip message-on-task, return
  // action: 'no-artifact'
}
\end{lstlisting}
\item TickResult union: \texttt{packages/holoscript-agent/src/types.ts}
line 47 declares \texttt{TickResult}; line 55 includes
\texttt{| 'no-artifact'} in the action union.
\item Tightening: \holoscript{} commit \texttt{435928977}
(2026-04-26 11:15 PT). \texttt{tools.ts} line 85 exports
\texttt{isProductiveBashCommand(cmd)} splitting bash commands
into \texttt{BASH\_READ\_ONLY\_PREFIXES} (cat, grep, ls, echo,
head, tail, wc, find, pwd, lake env, git status/log/diff/show)
and \texttt{BASH\_PRODUCTIVE\_PREFIXES} (lake build, lake clean,
lean, pnpm --filter, pnpm vitest, vitest run).
\end{itemize}

\textbf{Coverage by attack vector.}
\begin{center}\begin{tabular}{lll}
\toprule
Attack & Coverage & Reasoning \\
\midrule
A1 & High (by construction) & Pure-text response cannot pass the gate. \\
A2 & Medium (post-tightening) & Bounded by prefix-list completeness. \\
A3 & None & Gate is per-tick; impersonation is cross-agent. \\
A4 & None & Gate cannot observe brain-class drift. \\
A5 & None & Gate sees the commit succeed locally before upstream rejects. \\
\bottomrule
\end{tabular}\end{center}

\textbf{Measured efficacy.} Unit tests at
\texttt{packages/holoscript-agent/src/\_\_tests\_\_/runner.test.ts}
include three explicit no-artifact cases (lines 185, 214, 260):
pure-text response, read-only-bash-only, and empty-write\_file. All
three return \texttt{action: 'no-artifact'} and skip CAEL emission.
Test count for the runner suite: \textbf{113 \texttt{it()} blocks}
across 15 test files (14 in runner.test.ts, plus SHA extraction,
CAEL builder, audit log, ablation, brain loading, commit hooks, cost
guard, HoloMesh client, identity, integration, provisioning,
supervisor, and three-touch coverage). Aggregate suite:
\textbf{all passing} ($\geq 113/113$); pre-deploy snapshot was 109/111.

\subsection{D2: Fleet-Health Drift Detector}
\label{sec:d2}

\textbf{Mechanism.} A new global endpoint \texttt{GET
/api/holomesh/fleet/status} joins four signals per agent --- presence,
\cael{} record count, board claim age, distinct observed brain classes
--- and applies four drift rules:

\begin{enumerate}[topsep=2pt,itemsep=1pt]
\item D2.1 \texttt{cael\_no\_artifacts}: \cael{}-records-24h $> 0$
\textbf{and} done-entries-24h $= 0$.
\item D2.2 \texttt{stale\_claim}: claim-age-hours $> 6$ \textbf{and}
done-entries-24h $= 0$.
\item D2.3 \texttt{cael\_noisy\_no\_commits}: \cael{}-records-24h
$> 0$ \textbf{and} commits-on-done-24h $= 0$.
\item D2.4 \texttt{brain\_drift}: $\geq 2$ distinct brain\_class
values observed in 24h.
\end{enumerate}

A trust score per agent rolls these up: \texttt{trusted} (no flags)
$\rightarrow$ \texttt{warn} (1--2 flags) $\rightarrow$
\texttt{untrusted} ($\geq 3$ flags or D2.1 sustained $\geq 24\text{h}$).

\textbf{Citation.} \holoscript{} commit \texttt{95e3d898d}
(2026-04-26 10:58 PT).
\texttt{packages/mcp-server/src/holomesh/routes/core-routes.ts}
lines 758--997 define the rule set; the route handler joins
presence, \texttt{agentAuditStore}, and \texttt{team.taskBoard} into
the composite payload.

\textbf{Coverage by attack vector.}
\begin{center}\begin{tabular}{lll}
\toprule
Attack & Coverage & Reasoning \\
\midrule
A1 & High (retrospective) & D2.1 fires within 24h. \\
A2 & Medium (retrospective) & D2.3 catches when commits absent. \\
A3 & Low & DM-queue mismatch is not part of D2. \\
A4 & High (after 24h) & D2.4 is the canonical brain-drift catch. \\
A5 & Low & Detected only when \emph{all} commits fail; partial-fail invisible. \\
\bottomrule
\end{tabular}\end{center}

\textbf{Measured efficacy.} D2 is live in production. Over the 7-day
window post-deploy (2026-04-26 to 2026-05-03): 0 adversarial fires
(0 D2.1, 2 D2.2 operational stale-claim flags, 0 D2.3, 0 D2.4). Fleet
dormant during window; D2 would fire if conditions were met. Baseline
comparison: mw02 went undetected for 27.5h pre-D2 (2026-04-26). Backlog
test count: fleet-health route tests integrated into core-routes.ts; full
holomesh-suite count is $\geq 172$.

\subsection{D3: \cael{} \trustepoch{} Tag}
\label{sec:d3}

\textbf{Mechanism.} Every \cael{} audit record produced by a
post-W.107 runtime carries a tag
\texttt{trust\_epoch: 'post-w107'}. Records produced before the
2026-04-26 deploy carry no tag and are implicitly
\texttt{'pre-w107-untrusted'}. Downstream evidence pipelines (Paper
25 fleet gate clock; D2's drift detector; future paper-claim
verification scripts) filter by \trustepoch{} before treating CAEL as
evidence.

\textbf{Citation.} \holoscript{} commit \texttt{435928977}
(2026-04-26 11:15 PT).
\texttt{packages/holoscript-agent/src/cael-builder.ts}:
\begin{itemize}[topsep=2pt,itemsep=1pt]
\item Line 39--48: \texttt{trust\_epoch?: 'post-w107'} added to
\texttt{CaelAuditRecord} with documentation comment naming the
mw02 cohort.
\item Line 120: \texttt{buildCaelRecord} sets the tag on every new
record.
\end{itemize}

\textbf{Coverage by attack vector.} D3 is not an attack-blocker; it
is an \emph{evidence filter}. Its role is to ensure that the corpus
used by downstream agents is the post-defense corpus, not a
contaminated mixture. Coverage applies to evidence consumers
(Paper~25 gate clock, audit-log queries, paper-claim verification)
rather than to the attacks themselves.

\textbf{Measured efficacy.} Near-zero post-w107 records emitted during
the 7-day window (fleet dormant). The \texttt{trust\_epoch} tag is live
in production and every new CAEL record carries it. Excluding pre-W.107
records from D2's drift evaluation: not yet measured (fleet dormancy
produced insufficient post-w107 corpus to evaluate false-positive rate).

% =============================================================================
\section{Measurement Methodology}
\label{sec:methodology}
% =============================================================================

We compute attack-success rates against the corpus, not against the
routing scalar. The unit of measurement is \emph{a fabricated audit
record that survives in the corpus for $N$ wall-clock hours without
flag}. This differs from prior reputation-attack measurement
(Paper~1's $T(s)$-inflation methodology) because the corpus has no
decay --- a record persists until explicitly retired.

\textbf{Calibrated baseline.} The mw02 incident is a single trial
with attack-success $= 1$, observation window $= 27.5\text{h}$
(claim-time to defense-deploy). The 27.5h is a \emph{lower bound} on
how long A1 would have continued absent the fleet quality audit
that surfaced it; the true fundamental detection time absent any
defense is $\infty$ (no in-corpus signal would have crossed the
threshold).

\textbf{Defense efficacy formula.} Per attack class $X$ and defense
$D$:
\[ \text{efficacy}(D, X) = 1 - \frac{\text{success\_rate}(X \mid D)}
{\text{success\_rate}(X \mid \emptyset)}. \]
For $X = A1$ and $D = D1$: efficacy $= 1$ by construction (the gate
makes A1 success impossible at the runtime). For other cells:
the T7 snapshot chain records an honest-negative window: 0 adversarial
D2 fires, 2 operational stale-claim flags, and no adversarial harness
traffic, so rate estimates remain bounded rather than statistical.

\textbf{Post-defense baseline.} We ran the runner unit suite
post-deploy. Out of the runner test cases identified by grep, the
three explicit no-artifact assertions all pass on the gate
(\S\ref{sec:d1}). Aggregate \texttt{vitest} results for the
\texttt{@holoscript/agent} package: the D5 checkpoint records 4/4
no-artifact runner assertions passing; the full package suite was not
rerun in the snapshot chain, so we preserve the pre-deploy $109/111$
runner-suite figure as historical context rather than reporting it as
a fresh T7 pass rate.

% =============================================================================
\section{Results}
\label{sec:results}
% =============================================================================

\begin{table*}[t]
\centering
\small
\measuredFrom{research/paper-21-fleet-snapshots/2026-04-27-T0-baseline.json}%
\measuredFrom{research/paper-21-fleet-snapshots/2026-05-01-D5-checkpoint.json}%
\measuredFrom{research/paper-21-fleet-snapshots/2026-05-03-T7-final.json}%
\measuredFrom{scripts/lib/trust-epoch.mjs}%
\measuredFrom{scripts/paper-21-capture-T7.mjs}%
\measuredFrom{scripts/paper-21-extract-measurements.mjs}%
\measuredFrom{scripts/extract-sesl-pairs.mjs}%
\begin{tabular}{lcccccc}
\toprule
\textbf{Attack} & \textbf{Tier} & \textbf{Pre-defense} & \textbf{Defense covering it} & \textbf{Post-defense} & \textbf{Defense gap} \\
\midrule
A1: pure-monologue hallucination & T3 & $1.0$ over 27.5h (mw02) & D1 (by construction) + D2 (retrospective) & $0$ at runtime; 0 adversarial corpus fires (T7) & none \\
A2: trivial-bash bypass & T1/T3 & unmeasured & D1 (post-tightening) & bounded & prefix-list completeness \\
A3: recipient impersonation & T2 & unmeasured & operational mitigation only & unmeasured & no in-protocol defense \\
A4: brain-identity drift & T1+T3 & 1 instance (mw02) & D2.4 (retrospective) & 0 fires & $<24\text{h}$ window invisible \\
A5: credential shadow & T1 + env & 1 instance (W.075) & D2.3 (full-fail only) & 0 fires & partial-fail invisible \\
\bottomrule
\end{tabular}
\caption{Attack-success rates pre- and post-defense.
``Pre-defense'' columns measure success rate against the corpus
absent the cited defenses; the mw02 incident is the single
calibrated baseline available at submission time. The 7-day
post-deploy measurement window closed on 2026-05-03. A2--A5 cells
report honest negative results: defenses are deployed and live,
but fleet dormancy during the window prevented accumulation of
adversarial observations. See \S\ref{sec:methodology}.}
\label{tab:results}
\end{table*}

Table~\ref{tab:results} summarizes the measurement state at
submission time. The single calibrated cell (A1 pre-defense) is from
the mw02 incident; the remaining post-defense cells are honest-negative
or bounded by the dormant-fleet T7 snapshot chain.

The asymmetry between A1 and A4 is informative. A1 is a runtime
contract violation; D1 closes it by construction. A4 is a fleet-
configuration violation; D2.4 closes it only after the 24-hour
window the rule observes. The gap is the cost of the runtime not
having a per-tick brain-attestation check; we discuss in
\S\ref{sec:limitations} why we did not implement one.

% =============================================================================
\section{Limitations}
\label{sec:limitations}
% =============================================================================

\paragraph{A2 prefix-list maintenance.} D1's bash-side coverage
depends on \texttt{BASH\_PRODUCTIVE\_PREFIXES}. New evasions
(e.g.\ \texttt{npm run build} where \texttt{npm} is not in the
prefix list, or productive bash through unusual prefixes such as
\texttt{node script.js}) require list updates. We did not implement
a denylist (``these prefixes are read-only'') because the security
property we want is positive: the call must demonstrably advance
the work, not merely fail to be a no-op.

\paragraph{A3 in-protocol defense gap.} The current mitigation is
operational --- senders who use \texttt{/room send} for actionable
DMs include the literal phrase ``do not implement'' in the body, and
verify task state after delivery. An in-protocol defense would
require the subagent that handles \texttt{/room send} to be
authenticated as the \emph{sender}'s extension, with a hard prohibition
on executing tools in the recipient's role. We did not ship this in
the 2026-04-26 deploy because the change requires propagating an
authentication context into the skill's forked execution surface, which
is a multi-session refactor.

\paragraph{A4 sub-24h window.} D2.4 requires accumulating $\geq 2$
distinct brain\_class observations in a 24-hour window. An adversary
who completes their attack in $<24\text{h}$, then reverts the brain
configuration before the next heartbeat, evades D2.4. A per-tick
brain-attestation signature on every \cael{} record would close this
gap; we did not implement it because the attestation surface
overlaps with the seat-wallets signing protocol (workspace memory
entry S.SIGN) and a clean cutover requires the founder-attestation
flow that landed at 2026-04-25.

\paragraph{A5 partial-fail invisibility.} D2.3 requires
commits-on-done-24h $= 0$. The dangerous case is the agent that
\emph{usually} pushes successfully but silently fails on a subset of
pushes due to a credential shadow. The corpus accumulates records
of ``commit done'' for the failing subset that the upstream remote
will never confirm. The straightforward fix is a per-commit verify
step (resolve commit hash on the upstream remote within 60s of the
local commit); we did not ship it because the verify cost
($O(\text{N commits})$ per agent per day) exceeds the cost we are
willing to pay before the partial-fail rate is measured.

\paragraph{Single calibrated baseline.} The mw02 incident is one
trial. A1's pre-defense success rate of 1.0 is point-estimated, not
distribution-estimated. Production replay against a sandboxed
adversarial agent (per the existing
\texttt{research/paper-21-ati/evaluation-plan.md} testbed spec) is
required to bound the variance. This is scheduled for the post-deploy
measurement window.

\paragraph{Out of scope.} We do not address (1) the
\texttt{UISessionRecorder} attack surface scoped to Paper 24; (2)
formula-level attacks on Paper 1's $T(s)$ scalar (covered by the
companion \texttt{paper-21-ati/threat-model.md} document, with its
own five attack classes — Whitewasher, Sybil, Slow Poisoner, Score
Manipulator, Eclipse — and four defenses including Indistinguishable
Canary Probing); (3) cryptographic side channels.

% =============================================================================
\section{Related Work}
\label{sec:related}
% =============================================================================

\textbf{Reputation-system attacks.} The Sybil attack
\cite{TODO-find-citation-douceur-2002-sybil} and slow-poisoning
attacks on collaborative-filtering reputation systems
\cite{TODO-find-citation-recsys-poisoning} are the classical
foundations. The companion document
\texttt{paper-21-ati/threat-model.md} addresses these directly at
the trust-formula layer; this paper's attacks are at a strictly
higher layer (the runtime contract that produces the inputs the
formula consumes).

\textbf{LLM hallucination as a security property.} Recent work on
the security implications of LLM hallucinations
\cite{TODO-find-citation-llm-hallucination-security-2024} treats
hallucinations as a quality-of-service problem. Our framing is
different: in an autonomous mesh, an LLM's hallucinated final-text
is a write to a shared trust corpus, and the corpus is the security
boundary. The hallucination becomes \emph{the attack}, regardless of
the LLM's intent (T3 in our taxonomy).

\textbf{Prompt injection.} Greshake et al.\
\cite{TODO-find-citation-greshake-prompt-injection-2023} document
indirect prompt injection on LLM agents. A3 (recipient
impersonation) is an instance of indirect prompt injection on the
multi-agent surface --- the injected content is a message body
treated as instructions by the recipient-routing subagent.

\textbf{MCP supply-chain.} Hou et al.\
\cite{TODO-find-citation-hou-mcp-tool-poisoning} demonstrate
MCP-tool poisoning, where the tool annotation is the injection
vector. A5 (credential shadow) is a supply-chain attack on the
\texttt{.env}-injection path that initializes the agent's authentication
state; the attack vector is environment configuration rather than tool
annotation.

\textbf{Agent-mesh security.} \todo{find-citation: prior work on
agent-mesh security at USENIX Security 2023--2025 is sparse; the
closest comparison classes are P2P-network attacks and reputation
systems --- search engines: Google Scholar, USENIX Security
proceedings, IEEE S\&P proceedings 2023--2025}.

% =============================================================================
\section{Conclusion and Reproducibility}
\label{sec:conclusion}
% =============================================================================

We presented a formal threat model for runtime-contract attacks on
autonomous agent meshes, five attack vectors instantiated against
the \holoscript{} \holomesh{} fleet, and three deployed defenses.
The calibrated baseline is the mw02 incident of 2026-04-26 (A1
attack-success $= 1$ over 27.5h pre-defense). Post-defense, A1 is
blocked by construction at the runtime; A2 is blocked subject to
prefix-list completeness; A3, A4, and A5 have known gaps that we
flagged rather than concealed.

\paragraph{Reproducibility.} All defenses are open-sourced in the
\holoscript{} repository.
Cited commits: \texttt{4aab897ad}, \texttt{95e3d898d},
\texttt{435928977} (all on \texttt{main} as of 2026-04-26).
The companion brain-identity audit script lives at
\texttt{scripts/mesh-deploy/audit-brain-identity.mjs}. The audit task
that surfaced mw02 is workspace board task
\texttt{task\_1777223928419\_dmm5} (filed 2026-04-26 during the
fleet quality audit).

The complementary trust-formula attack model (Whitewasher, Sybil,
Slow Poisoner, Score Manipulator, Eclipse) is documented at
\texttt{research/paper-21-ati/threat-model.md} with its own three
defenses (exponential decay, cross-mesh anchoring, output diversity)
plus a fourth (indistinguishable canary probing) that closes a
formerly-residual gap on Slow Poisoner. Together the two documents
constitute Phase 1 of the Paper 21 program.

% =============================================================================
\section*{Acknowledgments}
% =============================================================================

This work was discovered by the \holoscript{} fleet's quality audit
on 2026-04-26 and the agents that participated in the post-incident
defense deploy. The companion threat model
(\texttt{paper-21-ati/threat-model.md}) was authored by the
\texttt{security-auditor-brain} on 2026-04-24.

% =============================================================================
\bibliographystyle{plain}
% =============================================================================

  % ═══════════════════════════════════════════════════════════════════════════
  % Bibliography — migrated 2026-05-03 from inline \begin{thebibliography} block
  % to shared research/holoscript.bib per founder /founder ruling
  % (research/paper-audit-matrix.md §2026-04-24 Refresh — Founder Ruling).
  %
  % Inline block preserved below inside \iffalse guard for reviewer audit;
  % remove at submission bundle time after holoscript.bib has been validated
  % to resolve every \cite key this paper uses.
  \bibliographystyle{plain}
  \bibliography{holoscript}

  \iffalse
  % === LEGACY INLINE BIBLIOGRAPHY (pre-2026-05-03 migration) ===
  % Kept behind \iffalse for audit trail only; bibtex ignores it.
  % Remove entirely at submission-bundle time after:
  %   pdflatex paper-21-adversarial-trust-injection-usenix && bibtex paper-21-adversarial-trust-injection-usenix
  % resolves 0 missing keys against research/holoscript.bib.
  \begin{thebibliography}{99}
\bibitem{mcp-spec}
Anthropic and Linux Foundation AI \& Data, \emph{Model Context
Protocol Specification}, 2024--2026.
\href{https://modelcontextprotocol.io}{modelcontextprotocol.io}.

\bibitem{TODO-self-cite-paper1}
Krzywoszyja, J. \emph{Trust by Replay: Hash-Verified Simulation
Results for Model Context Protocol Tool Use}. USENIX Security 2026
(submission). \texttt{research/paper-1-mcp-trust-usenix.tex}.

\bibitem{TODO-find-citation-douceur-2002-sybil}
\todo{cite: Douceur, J. R. ``The Sybil Attack.'' IPTPS 2002.}

\bibitem{TODO-find-citation-recsys-poisoning}
\todo{cite: survey paper on poisoning attacks against reputation
and recommender systems --- search ACM CCS / IEEE S\&P 2018--2024.}

\bibitem{TODO-find-citation-llm-hallucination-security-2024}
\todo{cite: 2024--2025 paper treating LLM hallucinations as a
security property --- candidate venues USENIX Security, IEEE S\&P,
NeurIPS Trustworthy ML.}

\bibitem{TODO-find-citation-greshake-prompt-injection-2023}
\todo{cite: Greshake et al., ``Not what you've signed up for:
Compromising Real-World LLM-Integrated Applications with Indirect
Prompt Injection.'' AISec 2023.}

\bibitem{TODO-find-citation-hou-mcp-tool-poisoning}
\todo{cite: Hou et al., MCP tool-poisoning paper, 2025.}

\bibitem{TODO-find-citation-agent-mesh-research-workloads}
\todo{cite: 2024--2026 paper documenting autonomous agent meshes on
production research workloads --- candidate sources include
multi-agent CHI/CSCW papers and AAMAS 2024--2026.}
\end{thebibliography}
  \fi

\end{document}
